\documentclass[12pt]{article}
\usepackage{graphicx} 
\usepackage{amsmath}
\usepackage[margin=1in]{geometry}
\usepackage{fancyhdr}
\usepackage{enumerate}
\usepackage[shortlabels]{enumitem}
\usepackage{amssymb}
\usepackage{amsthm}
\usepackage{cases}

%%%%%%%%%%%%%%%%%%%%%%%%%%%%%%%%%%%%%%%%%%%%%%%%%%%%%%%%%%%%%%%%%%%%%%
% FANCYHDR SETUP (For pages 2 and onwards)
%%%%%%%%%%%%%%%%%%%%%%%%%%%%%%%%%%%%%%%%%%%%%%%%%%%%%%%%%%%%%%%%%%%%%%
\pagestyle{fancy}
\fancyhead[l]{Giovanni Antonini}
\fancyhead[c]{Chapter 1: Vector Spaces}
\fancyfoot[c]{\thepage}
\renewcommand{\headrulewidth}{0.2pt} 
\setlength{\headheight}{15pt}        

%%%%%%%%%%%%%%%%%%%%%%%%%%%%%%%%%%%%%%%%%%%%%%%%%%%%%%%%%%%%%%%%%%%%%%
% PROBLEM ENVIRONMENT BOXES
%%%%%%%%%%%%%%%%%%%%%%%%%%%%%%%%%%%%%%%%%%%%%%%%%%%%%%%%%%%%%%%%%%%%%%
\usepackage{tcolorbox}
\tcbuselibrary{theorems, breakable, skins}
\newtcbtheorem{prob}% environment name
              {Problem}% Title text
  {enhanced, % tcolorbox styles
  attach boxed title to top left={xshift = 4mm, yshift=-2mm},
  colback=blue!5, colframe=black, colbacktitle=blue!3, coltitle=black,
  boxed title style={size=small,colframe=gray},
  fonttitle=\bfseries,
  separator sign none
  }%
  {} 
\newenvironment{problem}[1]{\begin{prob*}{#1}{}}{\end{prob*}}

%%%%%%%%%%%%%%%%%%%%%%%%%%%%%%%%%%%%%%%%%%%%%%%%%%%%%%%%%%%%%%%%%%%%%%
% SHORTCUT COMMANDS
%%%%%%%%%%%%%%%%%%%%%%%%%%%%%%%%%%%%%%%%%%%%%%%%%%%%%%%%%%%%%%%%%%%%%%
\newcommand{\F}{\mathbb{F}}
\newcommand{\R}{\mathbb{R}}
\newcommand{\C}{\mathbb{C}}
\begin{document}

% ====================================================================
% FIRST PAGE TITLE & CHAPTER BLOCK
% ====================================================================
\begin{center}
    \vspace*{1em}
    {\large \emph{Linear Algebra Done Right}, by Sheldon Axler} \\[0.5em]
    {\LARGE \bfseries Chapter 1: Vector Spaces} \\[0.3em]
    \vspace{1.5em}
\end{center}

% This single command tells LaTeX to hide the running headers on THIS page only:
\thispagestyle{plain} 

% ====================================================================
% CONTENT
% ====================================================================

% ====================================================================
% SECTION 1A: R^n and C^n
% ====================================================================
\section*{Section 1A: $\R^n$ and $\C^n$}

For problems 1-8, let $\alpha =a + bi$, $\beta = c+di$, and $\lambda = e+fi$ 

\begin{problem}{1}
Show that $\alpha + \beta =\beta + \alpha$ $\forall\alpha,\beta\in \mathbb{C}$
\end{problem}
\begin{proof}
    \begin{align*}
        \alpha + \beta &= (a+c)+(b+d)i \\
        &=(c+a)+(d+b)i \\
        &=\beta + \alpha
    \end{align*}
    Equalities 1 and 3 hold because of the definition of addition on $\mathbb{C}$ and equality 2 holds because of the commutativity of addition on $\mathbb{R}$.
\end{proof}

\begin{problem}{2}
Show that $(\alpha + \beta) + \lambda = \alpha + (\beta + \lambda)$ $\forall \alpha, \beta, \lambda \in \mathbb{C}$
\end{problem}
\begin{proof}
    \begin{align*}
        (\alpha + \beta) + \lambda &=[(a+c)+(b+d)i]+(e+fi)\\
        &=[(a+c)+e]+[(b+d)+f]i\\
        &=[a+(c+e)]+[b+(d+f)]i\\
        &=(a+bi)+[(c+e)+(d+f)i]\\
        &=\alpha + (\beta + \lambda)
    \end{align*}
    Equalities 1, 2, 4, and 5 hold by the definition of addition on $\mathbb{C}$. Equality 3 holds by the associative property of addition on $\mathbb{R}$.
\end{proof}

\begin{problem}{3}
Show that $(\alpha\beta)\lambda =\alpha(\beta \lambda)$ $\forall \alpha,\beta,\lambda \in \mathbb{C}$
\end{problem}
\begin{proof}
    \begin{align*}
        (\alpha\beta)\lambda &=[(ac-bd)+(ad+bc)i](e+fi)\\
        &= [(ac-bd)e-(ad+bc)f]+[(ac-bd)f+(ad+bc)e]i\\
        &= (ace-bde-adf-bcf) +(acf-bdf+ade+bce)i\\
        &= [a(ce-df)-b(cf+de)]+[a(cf+de)+b(ce-df)]i\\
        &= (a+bi)[(ce-df)+(cf+de)i]\\
        &= \alpha(\beta\lambda) 
    \end{align*}
    Equalities 1, 2, 5, and 6 hold by the definition of multiplication on $\mathbb{C}$. Equalities 3 and 4 hold by the commutative property of addition on $\mathbb{R}$ and the distributive property on $\mathbb{R}$.
\end{proof}

\begin{problem}{4}
Show that $\lambda(\alpha + \beta) = \lambda\alpha + \lambda\beta$ $\forall \alpha, \beta, \lambda \in \mathbb{C}$
\end{problem}
\begin{proof}
    \begin{align*}
        \lambda(\alpha + \beta) &= (e+fi)[(a+c)+(b+d)i]\\
        &= [e(a+c)-f(b+d)]+[e(b+d)+f(a+c)]i\\
        &=[(ea-fb)+(ec-fd)]+[(eb+fa)+(ed+fc)]i\\
        &=[(ea-fb)+(eb+fa)i]+[(ec-fd)+(ed+fc)i]\\
        &=\lambda\alpha +\lambda\beta
    \end{align*}
    Equalities 1, 2, and 5 hold true by the definition of multiplication on $\mathbb{C}$.
    Equality 3 holds true because of the commutativity of addition on $\mathbb{R}$ and by the distributive property on $\mathbb{R}$. Finally, equality 4 holds true by the definition of addition on $\mathbb{C}$.
\end{proof}

\begin{problem}{5}
Show that for every $\alpha \in \mathbb{C}$, there exists a unique $\beta \in \mathbb{C}$ such that $\alpha + \beta =0$
\end{problem}
\begin{proof}
    Let $\alpha=a + bi$ and $\beta = c+di$. 
    \[\alpha + \beta =0\]
    can be rewritten as
    \[(a+bi)+(c+di)=0\]
    and then as
    \[(a+c)+(b+d)i=0\]
    by the definition of addition on $\mathbb{C}$. Then $c$ must be the unique additive inverse of $a$ in $\mathbb{R}$ and $d$ must be the unique additive inverse of $b$ in $\mathbb{R}$. Recall there exists a bijection
    \[f: \mathbb{R}^2\to \mathbb{C}\]
    therefore for every $(r_1,r_2)\in \mathbb{R}^2$ there exists one unique $c\in \mathbb{C}$. Therefore, the pair $(c,d)\in \mathbb{R}$ is mapped to the unique number $c+di=\beta\in\mathbb{C}$.
\end{proof}

\begin{problem}{6}
Show that for every $\alpha \in \mathbb{C}$ with $\alpha \ne 0$, there exists a unique $\beta \in \mathbb{C}$ such that $\alpha \beta=1$
\end{problem}
\begin{proof}
    Let $\alpha = a + b$ and $\beta = b+di$. Then
    \[\alpha\beta =1 \implies (a+bi)(c+di)=1 \implies (ac-bd)+(ad+bc)i=1\]
    Recall that for $1\in \mathbb{R}\subset \mathbb{C}$, $1=1+0i$. So, 
    \begin{align}
        ac-bd=1 \\
        ad+bc=0
    \end{align}
    Therefore we can solve for $c$ and $d$ in terms of $a$ and $b$. Rewriting (1)
    \begin{align}
        c=\frac{1+bd}{a}
    \end{align}

    Now we can substitute $c$ into (2):
    \begin{align*}
        ad+b(\frac{1+bd}{a})&=0\\
        ad+\frac{b}{a}+\frac{b^2d}{a}&=0\\
        d(a+\frac{b^2}{a})+\frac{b}{a}&=0\\
    \end{align*}
    And solving for $d$:
    \begin{align*}
        d&=-\frac{b}{a} \left(\frac{1}{a+\frac{b^2}
        {a}}\right)\\ 
        &d= \frac{-b}{a^2+b^2}
    \end{align*}
    Lastly we can substitute $d$ into (3) and we get:
    \begin{align}
        (c,d)=\left(\frac{a}{a^2+b^2},\frac{-b}{a^2+b^2}\right)
    \end{align}
    From this we determine that there exists $f:\mathbb{R}^2\to \mathbb{R}^2$ defined by
    \[f(x,y)=\left(\frac{x}{x^2+y^2},\frac{-y}{x^2+y^2}\right)\]
    which maps $(a,b)$ to $(c,d)$. We must now determine if this function is bijective in order to show that for every $(a,b)\in \mathbb{R}^2$, there exists a unique pair $(c,d)\in\mathbb{R}^2$ which then is mapped to a unique complex number $c+di\in \mathbb{C}$.\\
    Recall that a function is bijective if the inverse of such function is the function itself. In other words:
    \[f \circ f = id_X \quad \text{or} \quad f(f(x,y))=(x,y)\]
    where 
    \[X=\{(x,y)\in\mathbb{R}^2\mid x,y)\ne (0,0)\}\]
    Let $(x',y')=f(x,y)$, so
    \[x'=\frac{x}{x^2+y^2} \quad \text{and} \quad y'=\frac{-y}{x^2+y^2}\]
    Then 
    \[f(x',y')=\left(\frac{x'}{x'^2+y'^2},\frac{-y'}{x'^2+y'^2}\right)\]
    Rewriting the denominators separately
    \[x'^2+y'^2=\left(\frac{x}{x^2+y^2}\right)^2+\left(\frac{-y}{x^2+y^2} \right)^2=\frac{1}{x^2+y^2}\]
    Therefore we can simplify
    \[f(x',y')=\left(\frac{x'}{\frac{1}{x^2+y^2}},\frac{-y'}{\frac{1}{x^2+y^2}} \right)=(x,y)\]
    Since $f(f(x,y))=(x,y)$, we can say that $f$ is a bijection. Therefore, for $(a,b)\in \mathbb{R}^2$, there exists a unique pair $(c,d)\in \mathbb{R}^2$. Then, there also exists a unique complex number $c+di=\beta \in \mathbb{C}$
\end{proof}

\begin{problem}{7}
Show that \[\frac{-1+\sqrt{3}i}{2}\] is a cube root of 1
\end{problem}
\begin{proof}
    We must show that $(\frac{-1+\sqrt{3}i}{2})^3=1$:
    \begin{align*}
        \left(\frac{-1+\sqrt{3}i}{2}\right)^3&=(-1+\sqrt{3}i)^3(1/2)^3\\
        &=(-1+\sqrt{3}i)(-1+\sqrt{3}i)(-1+\sqrt{3}i)(1/8)\\
        &=[(-2)+(-2\sqrt{3})](-1+\sqrt{3}i)(1/8)\\
        &=[(2-(-6)+(-2\sqrt{3}+2\sqrt{3})i](1/8)\\
        &=(8+0i)(1/8)\\
        &=8(1/8)\\
        &=1
    \end{align*}
\end{proof}

\begin{problem}{8}
Find two distinct square roots of $i$
\end{problem}
\begin{proof}
    We must find some $c\in \mathbb{C}$ such that 
    \[c^2=i\]
    Let $c=a+bi$, then
    \[c^2=(a+bi)(a+bi)=(a^2-b^2)+(2ab)i=i\]
    This means that we have the following
    \begin{align}
        2ab=1\\
        a^2-b^2=0
    \end{align}
    Rewriting (6)
    \[a=b\]
    then substituting into (5)
    \[2a^2=1\implies a=\pm \frac{1}{\sqrt{2}}\]
    Since for every $z\in\mathbb{C}$ there exists a number (an additive inverse) $-z\in\mathbb{C}$ such that
    \[z+(-z)=1\]
    then we have two numbers in $\mathbb{C}$ which are the square roots of $i$:
    \[c=\frac{1}{\sqrt{2}}+\frac{1}{\sqrt{2}}i \quad \text{and} \quad -c=-\frac{1}{\sqrt{2}}+-\frac{1}{\sqrt{2}}i\]
    where $-c$ is the additive inverse of $c\in \mathbb{C}$.  
\end{proof}

\begin{problem}{9}
Find $x\in \mathbb{R}^4$ such that \[(4,-3,1,7)+2x=(5,9,-6,8)\]
\end{problem}
\begin{proof}
    Let $x=(a,b,c,d)$. Then by the property of scalar multiplication in $\mathbb{R}^4$, and more generally $\mathbb{R}^n$
    \[2x=2(a,b,c,d)=(2a,2b,2c,2d)\]
    So
    \begin{align*}
            (4,-3,1,7)+2x&=(4,-3,1,7)+(2a,2b,2c,2d)\\
            &=(4+2a,-3+2b,1+2c,7+2d)\\
            &=(5,9,-6,8)
    \end{align*}
    where equality 2 holds by the definition of addition in $\mathbb{R}^4$.
    Then we have 4 equations to solve:
    \[4+2a=5\]
    \[-3+2b=9\]
    \[1+2c=-6\]
    \[7+2d=8\]
    Which produces the 4-tuple
    \[(a,b,c,d)=(\frac{1}{2},6,-\frac{7}{2},\frac{1}{2})\]
\end{proof}

\begin{problem}{10}
Explain why there does not exist $\lambda \in \mathbb{C}$ such that \[\lambda(2-3i,5+4i,-6+7i)=(12-5i,7+22i,-32-9i)\]
\end{problem}
\begin{proof}
    Let $\lambda =a+bi$. Then by the property of scalar multiplication in $\mathbb{C}^3$
    \[\lambda(2-3i,5+4i,-6+7i)=((a+bi)(2-3i),(a+bi)(5+4i),(a+bi)(-6+7i))\]
    which then by the definition of multiplication in $\mathbb{C}$ becomes
    \[((2a+3b)+(2b-3a)i,(5a-4b)+(4a+5b)i,(-6a-7b)+(7a-6b)i)\]
    Setting it equal to $(12-5i,7+22i,-32-9i)$ produces the following equalities:
    \[(2a+3b)+(2b-3a)i=12-5i\]
    \[(5a-4b)+(4a+5b)i=7+22i\]
    \[(-6a-7b)+(7a-6b)i=-32-9i\]
    Analyzing the first equality we can solve for $a$ and $b$
    \begin{align}
        2a+3b=12\\
        2b-3a=-5
    \end{align}
    Multiplying (7) by 3 and (8) by 2 we get
    \[6a+9b=36\]
    \[-6a+4b=-10\]
    Adding them together:
    \[13b=26\implies b=2\]
    Substituting back into (7)
    \[2a+3(2)=12\implies a=3\]
    So $\lambda =3+2i$. But if we check this pair for the third equality we will see that while the real portion holds
    \[-6(3)-7(2)=-32\]
    the imaginary portion does not
    \[(7(3)-6(2))i=9i\ne -9i\]
\end{proof}

\begin{problem}{11}
Show that $(x+y)+z=x+(y+z)$ $\forall x,y,z\in \mathbb{F}^n$
\end{problem}
\begin{proof}
    Recall that $\mathbb{F}$ can be either $\mathbb{R}$ or $\mathbb{C}$. 
    Let $x=(x_1,\ldots,x_n),y=(y_1,\ldots,y_n),$ and $z=(z_1,\ldots, z_n)$. Then
    \begin{align*}
        (x+y)+z&=((x_1,\ldots,x_n)+(y_1,\ldots,y_n))+(z_1,\ldots,z_n)\\
        &=((x_1+y_1),\ldots,(x_n+y_n))+(z_1,\ldots,z_n)\\
        &=((x_1+y_1)+z_1,\ldots, (x_n+y_n)+z_n)\\
        &=(x_1+(y_1+z_1),\ldots,x_n+(y_n+z_n))\\
        &=(x_1,\ldots,x_n)+(y_1+z_1,\ldots,y_n+z_n)\\
        &=x+(y+z)
    \end{align*}
    Equalities 1, 2, 3, 5, and 6 hold by the definition of vector addition in $\mathbb{F}^n$. Equality 4 holds by the associative property of addition on $\mathbb{F}$.
\end{proof}

\begin{problem}{12}
Show that $(ab)x=a(bx)$ $\forall x\in \mathbb{F}^n$ and $\forall a,b\in \mathbb{F}$
\end{problem}
\begin{proof}
    Let $x=(x_1,\ldots,x_n)$
    \begin{align*}
        (ab)x&=((ab)x_1,\ldots,(ab)x_n)\\
        &=(a(bx_1),\ldots,a(bx_n))\\
        &=a(bx_1,\ldots,bx_n)\\
        &=a(bx)
    \end{align*}
    Equality 1, 3, and 4 holds by the definition of scalar multiplication in $\mathbb{F}^n$. Equality 2 holds by the associative property of multiplication in $\mathbb{F}$.
\end{proof}

\begin{problem}{13}
Show that $1x=x$ $\forall x\in \mathbb{F}^n$
\end{problem}
\begin{proof}
    Let $x=(x_1,\ldots,x_n)$
    \begin{align*}
        1x&=1(x_1,\ldots,x_n)\\
        &=(1x_1,\ldots,1x_n)\\
        &=(x_1,\ldots,x_n)\\
        &=x
    \end{align*}
    Equality 2 holds true by the definition of scalar multiplication in $\mathbb{F}$. Equality 3 holds true since $1$ is the multiplicative identity in $\mathbb{F}$.
\end{proof}

\begin{problem}{14}
Show that $\lambda(x+y)=\lambda x+\lambda y$
\end{problem}
\begin{proof}
    Let $x=(x_1,\ldots ,x_n)$ and $y=(y_1,\ldots,y_n)$.
    \begin{align*}
        \lambda (x+y)&=\lambda((x_1,\ldots, x_n)+(y_1,\ldots,y_n))\\
        &=\lambda(x_1+y_1,\ldots,x_n+y_n)\\
        &=(\lambda (x_1+y_1),\ldots, \lambda(x_n+y_n))\\
        &=(\lambda x_1 + \lambda y_1,\ldots, \lambda x_n +\lambda y_n)\\
        &=(\lambda x_1,\ldots, \lambda x_n)+(\lambda y_1,\ldots, \lambda y_n)\\
        &=\lambda x + \lambda y
    \end{align*}
    Equalities 1,2, and 5 hold by the definition of addition in $\mathbb{F}^n$. Equalities 3 and 6 hold by the definition of scalar multiplication in $\mathbb{F}^n$. Equality 4 holds by the distributive property in $\mathbb{F}$. 
\end{proof}

\begin{problem}{15}
Show that $(a+b)x=ax+bx$ $\forall a,b\in \mathbb{F}$ and $\forall x\in \mathbb{F}^n$
\end{problem}
\begin{proof}
    Let $x=(x_1,\ldots,x_n)$.
    \begin{align*}
        (a+b)x&=(a+b)(x_1,\ldots,x_n)\\
        &=((a+b)x_1,\ldots,(a+b)x_n)\\
        &=(ax_1+bx_1,\ldots,ax_n+bx_n)\\
        &=(ax_1,\ldots, ax_n)+(bx_1,\ldots,bx_n)\\
        &=ax+bx
    \end{align*}
    Equalities 2 and 4 hold by the definition of scalar multiplication in $\mathbb{F}^n$. Equality 3 holds by the distributive property in $\mathbb{F}$.
\end{proof}

\newpage

% ====================================================================
% SECTION 1B: Definition of a Vector Space
% ====================================================================
{\large \bfseries Section 1B: Definition of a Vector Space} \\
We will define $V$ as a vector space over $\F$.

\begin{problem}{1}
Prove that $-(-v)=v$ for all $v \in V$
\end{problem}
\begin{proof}
We will assume $(-1)v=-v$ for every $v \in V$ as shown in the textbook under 1.32.
\begin{align*}
-(-v) &= (-1)(-v) && \text{by 1.32} \\
      &= (-1)((-1)v) && \text{by 1.32 inside the parenthesis} \\
      &= ((-1)(-1))v && \text{by associativity of scalar multiplication} \\
      &= (1)v        && \text{since } (-1)(-1) = 1 \text{ in } \F \\
      &= v           && \text{by the multiplicative identity of } V
\end{align*}
\end{proof}

\begin{problem}{2}
    Suppose $a\in \F, v\in V$, and $av=0$. Prove that $a=0$ or $v=0$
\end{problem}
\begin{proof}
    Note, we are relying on 1.30 and 1.31 from the textbook which assert the following:
    \[0v=0 \quad \forall v\in V\]
    and
    \[a0=0 \quad \forall a \in \F\]
    Now, we can write the statement we are trying to prove as
    \[av=0\implies (a=0 \lor v=0)\]
    which is logically equivalent to 
    \[\neg(av=0)\lor (a=0\lor v=0)\]
    For the sake of contradiction, we negate the statement 
    \[\neg(\neg(av=0)\lor (a=0 \lor v=0))\equiv (av=0)\land (a\ne0\land v\ne0)\]
    RAA: let $av=0$, $a\ne 0$, and $v \ne 0$. 
    Then 
    \[av=a^{-1}av=a^{-1}0=0\]
    but also
    \[av=a^{-1}av=1v=v\ne 0\]
    therefore contradiction. Thus, the original statement holds true.
\end{proof}
    \newpage
    \begin{problem}{3}
        Suppose $v,w,\in V$. Explain why there exists a unique $x\in V$ such that $v+3x=w$
    \end{problem}
    \begin{proof}
    We will start by showing the existence of $x$. 
    Let $x=3^{-1}(w+(-v))$
    \begin{align*}
        v+3(3^{-1}(w+(-v)))&=v+3(3^{-1}w+3^{-1}(-v))\\
        &=v+3(3^{-1}w)+3(3^{-1}(-v))\\
        &=v+(3(3^{-1}))w+(3(3^{-1}))(-v)\\
        &=v+1w+1(-v)\\
        &=v+(-v)+w\\
        &=0+w\\
        &=w
    \end{align*}
    Therefore, $x=3^{-1}(w+(-v))$ is a solution. Now we must show it is unique. Suppose $x'$ is also a solution. We will work under the assumption that $\forall a\in \F$, $-a=-1a$.
    \begin{align*}
        0&=w+(-w)\\
        &=(v+3x)+(-(v+3x'))\\
        &=(v+3x)+(-1(v+3x'))\\
        &=(v+3x)+(-1v+(-1)(3x'))\\
        &=(v+3x)+(-v+(-1(3))x')\\
        &=(v+3x)+(-v+(3(-1))x')\\
        &=v+(-v)+3x+3(-1x')\\
        &=0+3x+3(-x')\\
        &=3(x+(-x'))
    \end{align*}
    As shown in 1B problem 2, either $3=0$ or $(x+(-x'))=0$. Since $3\ne 0$, then
    \begin{align*}
        x+(-x')&=0\\
        x+(-x')+x'&=0+x'\\
        x+0&=x'+0\\
        x&=x'
    \end{align*}
    Thus, $x$ is unique. 
    \end{proof}

    \begin{problem}{4}
        The empty set is not a vector space. The empty set fails to satisfy only one of the requirements listed in the definition of a vector space (1.20). Which one?
    \end{problem}
    \begin{proof}
    If we rewrite the 6 properties in 1.20 with logical statements, we will see that they all, but one, have as their antecedent the existence of some vector(s) in the vector space $V$. For instance, the commutative property of addition:
    \[(\forall u,v\in V)(u,v\in V \implies u+v=v+u)\]
    Thus, if $V\subseteq\emptyset$, then the statement is vacuously true as no elements exist.
    However, the one dissimilar property is the additive identity as it does not depend on the existence of vectors within the vector space. It is a proposition that demands the existence of an element $0\in V$. Therefore, since $0\notin \emptyset$, then the empty set does not have an additive identity. This means the empty set is not a vector space.
    \end{proof}
        
    \begin{problem}{5}
        Show that in the definition of a vector space (1.20), the additive inverse condition can be replaced with the condition that
        \[0v=0 \quad \text{for all }v\in V\]
        Here the 0 on the left side is the number 0, and the 0 on the right side is the additive identity of $V$.
    \end{problem}
    \begin{proof}
    We must show that the properties of a vector space \textbf{with} the additive inverse property imply $0v=0$. Then we must show that the properties \textbf{without} the additive inverse property and \textbf{with} $0v=0$ imply the additive inverse property.\\
    $\implies$\\
    Suppose we have the usual properties of a vector space $V$. Then
    \begin{align*}
        0v&=(0+0)v\\
        &=0v+0v
    \end{align*}
    Then if we add the additive inverse of $0v$ to both sides we get
    \[0=0v\]
    $\Longleftarrow$\\
    Suppose $0v=0$ but we don't formally have the additive inverse property.
    \begin{align*}
        0v&=(1+(-1))v=1v+(-1)v=v+(-1)v=0
    \end{align*}
    Therefore, $(-1)v$ is the inverse of $v$ $\forall v\in V$.
    Note: we know $0=1+(-1)$ by the additive inverse property in $\F$.
    Therefore, the two conditions/properties can be interchanged. 
    \end{proof}
    \begin{problem}
        Let $\infty$ and $-\infty$ denote two distinct objects, neither of which is in $\R$. Define an addition and scalar multiplication on $\R \cup \{\infty,-\infty\}$ as you could guess from the notation. Specifically, the sum and product of two real numbers is as usual, and for $t\in \R$ define
        \[
        t\infty=\begin{cases}
            -\infty \quad &\text{if   } t<0,\\
            0  &\text{if   } t=0,\\
            \infty &\text{if   } t>0,
        \end{cases}
        \qquad t(-\infty)=\begin{cases}
            \infty \quad &\text{if   } t<0,\\
            0 &\text{if   } t=0,\\
            -\infty &\text{if   } t>0,
        \end{cases}
        \]
    and
    \[t+\infty = \infty + t = \infty + \infty = \infty\]
    \[t+(-\infty)=(-\infty)+t=(-\infty) +(-\infty)=-\infty\]
    \[\infty + (-\infty)=(-\infty)+\infty =0\]
    With these operations of addition and scalar multiplication, is $\R \cup \{\infty, -\infty\}$ a vector space over $\mathbb{R}$? Explain.
    \end{problem}
    \begin{proof}
    Immediately we can prove $\R \cup \{\infty,-\infty\}$ over $\R$ is not a vector space by checking for associativity of addition. 
    \[(\infty + \infty )+(-\infty)=\infty +(-\infty)=0\]
    \[\infty +(\infty +(-\infty))=\infty + 0=\infty\]
    Since associativity of addition does not hold $\forall a,b,c\in \R\cup\{\infty, -\infty\}$, then $\R \cup \{\infty, -\infty\}$ is not a vector space over $\R$.
    \end{proof}

    \begin{problem}{7}
        Suppose $S$ is a nonempty set. Let $V^S$ denote the set of functions from $S$ to $V$. Define a natural addition and scalar multiplication on $V^S$, and show that $V^S$ is a vector space with these definitions.
    \end{problem}
    Assuming $V$ is a vector space over $\F$.\\
    \noindent \textbf{Addition}: for $f,g\in V^S$, $x\in S$, and $f(x),g(x)\in V$, addition is defined as 
    \[(f+g)(x)=f(x)+g(x)\]
    \noindent \textbf{Scalar Multiplication}: for $f\in V^S$, $x\in S$, $\lambda \in \F$, and $f(x)\in V$, scalar multiplication is defined as
    \[(\lambda f)(x)=\lambda f(x)\]
    Now we will prove how $V^S$ inherits the properties of a vector space.\\
    \noindent \textbf{Commutativity}: for all $f,g\in V^S$ and $x\in S$, since $f(x),g(x)\in V$ and addition is commutative in $V$,
    \[(f+g)(x)=f(x)+g(x)=g(x)+f(x)=(g+f)(x)\]

    \noindent \textbf{Associativity}: for all $f,g,h\in V^S$, $x\in S$, since $f(x),g(x),h(x)\in V$ and addition is associative in $V$,
    \begin{align*}
        (f+g)(x)+h(x)&=[f(x)+g(x)]+h(x)\\
        &=f(x)+[g(x)+h(x)]\\
        &=f(x)+(g+h))(x)
    \end{align*}
    Now for associativity of scalar multiplication. Let $a,b\in \F$. Since scalar multiplication is associative in $V$,
    \[(a(bf))(x)=a[(bf)(x)]=a[bf(x)]=(ab)f(x)\]
    \noindent \textbf{Additive Identity}: first, define $0(x)\in V^S$, the additive identity of $V^S$
    \[0(x)=0\]
    for all $x\in S$.\\
    Now, let $f\in V^S$. Since $f(x)\in V$ and $0$ is the additive identity in $V$, 
    \[(f+0)(x)=f(x)+0(x)=f(x)+0 = f(x)\]
    \noindent \textbf{Additive Inverse}: for all $f\in V^S$, define $(-f)\in V^S$ as its inverse such that
    \[(-f)(x)=-f(x)\]
    for all $x\in S$. Since $-f(x)\in V$
    \[(f+(-f))(x)=f(x)+(-f(x))=0\]
    \noindent \textbf{Multiplicative Identity}: consider the scalar $1\in \F$. For all $f\in V^S$, $x\in S$, by definition of scalar multiplication in $V^S$ we have
    since $f(x)\in V$,
    \[(1f)(x)=1f(x)\]
    and since $f(x)\in V$
    \[1f(x)=f(x)\]
    \noindent \textbf{Distributive Property}: for all $f,g\in V^S$, $a\in \F$ and $x\in S$, by the definition of scalar multiplication in $V^S$ we have
    \[(a(f+g))(x)=a((f+g)(x))=a(f(x)+g(x))\]
    and since $f(x),g(x)\in V$
    \[a[f(x)+g(x)]=af(x)+ag(x)\]
    which is equivalent to
    \[(af)(x)+(ag)(x)\]
    Now, to distribute scalars over a vector, let $b\in \F$, by the definition of scalar multiplication we have
    \[((a+b)f)(x)=(a+b)f(x)\]
    and since $f(x)\in V$
    \[(a+b)f(x)=af(x)+bf(x)\]
    which is equivalent to 
    \[(af)(x)+(bf)(x)\]

    \begin{problem}{8}
        Suppose $V$ is a real vector space.
        \begin{itemize}
            \item The $\textit{complexification}$ of $V$, denoted by $V_C$, equals $V\times V$. An element of $V_C$ is an ordered pair $(u,v)$, where $u,v\in V$, but we write this as $u+iv$.
            \item Addition on $V_C$ is defined by
            \[(u_1+iv_1)+(u_2+iv_2)=(u_1+u_2)+i(v_1+v_2)\]
            for all $u_1,v_1,u_2,v_2\in V$
            \item Complex scalar multiplication on $V_C$ is defined by
            \[(a+bi)(u+iv)=(au-bv)+i(av+bu)\]
            for all $a,b\in \R$ and all $u,v\in V$.
        \end{itemize}
        Prove that with the definitions of addition and scalar multiplication as above, $V_C$ is a complex vector space.\\
        \textit{Think of $V$ as a subset of $V_C$ by identifying $u\in V$ with $u+i0$. The construction of $V_C$ from $V$ can then be thought of as generalizing the construction of $\C^n$ from $\R ^n$. }
    \end{problem}
    \begin{proof}
    Recall that we showed $\C$ had the properties of being commutative under addition, associative under addition and scalar multiplication, an additive and multiplicative identity, a unique additive inverse for every element, and being both left and right distributive. Since there exists an injective function
    \[f:V_C\to \C\]
    where \[f(u,v)=u+iv\in \C\]
    then we can treat $V_C$ as a subset of $\C$ or a set of complex numbers. With that, $V_C$ is naturally commutative and associative over addition. 
    \\Note, we are checking to see if $V_C$ is a complex vector space, therefore the scalar multiplication is dependent on $\C$. So, scalar multiplication is associative due to the associativity of multiplication on $\C$ (since $V_C\subset \C)$. 
    \\Then, since $V$ is a vector space, it naturally has an additive identity $0\in V$ and a multiplicative identity $1\in V$. Therefore, $V_C$ also has such elements as $V\subset V_C$.
    \\Next, for every $v\in V$, there exists $-v\in V$ such that $v+(-v)=0$. So, if we have $(a,b)\in V_C$, there exists $-a\in V$ and $-b\in V$. This in turn means that there exists $(-a,-b)\in V_C$. Since $(-a,-b)$ is the additive inverse of $(a,b)$
    \[(a+ib)+(-a+i(-b))=(a+(-a))+i(b+(-b))=0+i0=0\]
    then $V_C$ fulfills the additive inverse property.
    \\Lastly, recall the distributive property is dependent on scalars in $\C$. Since elements in $V_C$ are elements in $\C$, then this simplifies to the distributive property on $\C$ which holds true.
    
    \end{proof}
\end{document}